% Generated draft: requires mathematical and source review.
\documentclass[11pt]{article}
\usepackage[margin=25mm]{geometry}
\usepackage{fontspec}
\setmainfont{Latin Modern Roman}
\usepackage{amsmath,amssymb,mathtools,unicode-math}
\setmathfont{Latin Modern Math}
\usepackage{enumitem,needspace,xurl,hyperref}
\setlength{\parindent}{0pt}
\setlength{\parskip}{5pt}
\setlength{\emergencystretch}{3em}
\newcommand{\N}{\mathbb{N}}
\newcommand{\Z}{\mathbb{Z}}
\newcommand{\Q}{\mathbb{Q}}
\newcommand{\R}{\mathbb{R}}
\newcommand{\C}{\mathbb{C}}
\newcommand{\F}{\mathbb{F}}
\newcommand{\A}{\mathbb{A}}
\newcommand{\PP}{\mathbb{P}}
\newcommand{\E}{\mathbb{E}}
\newcommand{\eps}{\varepsilon}
\newcommand{\dd}{\,\mathrm{d}}
\begin{document}
\section*{UnsolvedMath ID: 800227}
\textbf{Source problem code:} catalogue entry
\subsection{Definitions and mathematical statement}
Let \(n\geq 1\) be an integer, and let \(\Omega\subset\C^n\) be a domain, meaning a connected open set, with smooth boundary, assumed pseudoconvex. Define
\[
H^\infty(\Omega)=\left\{f\in\mathcal O(\Omega):\sup_{z\in\Omega}|f(z)|<\infty\right\}.
\]
Write \(|V|\) for the support of the divisor \(V\) in \(\Omega\). The source assumes that \(V\) is bounded away from the \emph{Silov (Shilov) boundary} in \(\partial\Omega\). If that boundary is denoted by \(S\), the separation condition is
\[
\operatorname{dist}(|V|,S)
=\inf\{\|z-\zeta\|:z\in|V|,\ \zeta\in S\}>0.
\]
The function algebra determining \(S\) is not specified in the source, so this hypothesis is not yet fully determined. The wording ``smoothly bounded'' also leaves unclear whether boundedness of \(\Omega\), in addition to smoothness of its boundary, is required.

With vanishing interpreted as vanishing on the support, the literal requested conclusion is
\[
\exists f\in H^\infty(\Omega)\quad
\forall z\in|V|,\quad f(z)=0.
\]
The source does not explicitly require \(f\not\equiv0\). As written, this conclusion is satisfied by the zero function. The apparent nontrivial intended conclusion adds \(f\not\equiv0\), but that is an interpretation requiring confirmation.

The source also does not specify whether \(V\) is effective or whether its multiplicities must be respected. If \(V=\sum_j m_jY_j\) is an effective analytic divisor, where the \(Y_j\) are irreducible analytic hypersurfaces, the sum is locally finite, and \(m_j\) are positive integers, vanishing to the prescribed multiplicities would mean
\[
\operatorname{ord}_{Y_j}(f)\geq m_j\qquad\text{for every }j,
\]
for a nonzero holomorphic \(f\). This is stronger than merely vanishing on \(|V|\); the supplied wording does not distinguish these conclusions.

The source reports that the result is known for the polydisc and attributes that result to Rudin. This qualification is retained as a supplied claim; its precise formulation and reference are not supplied.
\subsection{Short English statement}
The source asks whether a divisor separated from the Shilov boundary of a smoothly bounded pseudoconvex domain is contained in the zero set of a bounded holomorphic function. A nonzero function is apparently intended but is not explicitly required. The relevant function algebra, boundedness assumption, and treatment of divisor multiplicities require clarification. The supplied record labels the problem open and attributes a polydisc result to Rudin; these are source claims, not independently verified assessments.
\subsection{Sources}
\begin{itemize}
\item Supplied local UnsolvedMath registry record: permanent ID \(800227\), exact source problem code \texttt{catalogue entry}, title ``Several Complex Variables Problem List --- Problem 23.'' The registry lists Buzzard, Fornaess, and Krantz as proposers, with proposed year 1995, and labels the record open.
\item Buzzard, Fornaess, and Krantz, \emph{The Several Complex Variables Problem List} (1995), Problem 23. Supplied source URL: \url{https://arxiv.org/abs/math/9509202}. The supplied provenance records access on 2026-10-16 and extraction by \texttt{arxiv\_latex}; these provenance details are not independently verified here.
\item The supplied problem text is signed S. Krantz, with contact printed as \texttt{sk @math.wustl.edu}.
\item The problem text attributes the polydisc case to Rudin, but supplies no bibliographic citation for that result.
\end{itemize}
\end{document}
